\chapter{The Film as a Family of Realizations}\label{ch:family}

Part~I established what the screen can carry. This part asks what it should carry. Two unrelated films of equal length can be interleaved, and the result is two films that happen to share a room. The proposal of this book is something else: a single work that exists as several related realizations, designed together and meant to be seen together. This chapter defines that object and distinguishes it from its weaker neighbors.

\section{Realization and family}

A \emph{realization} is a complete audiovisual experience, one film a viewer can follow from beginning to end. A \emph{variant family} is a set of realizations
\begin{equation}
\Fam=\{F_1,F_2,\ldots,F_m\}
\end{equation}
designed for synchronized presentation. Each member is whole. None is a supplement, a commentary track, or a set of optional scenes attached to another.

What makes a set of realizations a family rather than a program of short films is a shared \emph{temporal scaffold}: a common clock against which every realization is laid out, and a set of moments at which the realizations are required to correspond. The scaffold is what lets several people in one room remain in the same occasion while watching different films, and it is what Part~I's adaptive gating exploits when the realizations coincide on screen.

\section{Degrees of relatedness}

Realizations can be related more or less closely, and the degree determines what the family can do.

At the weakest level, the realizations merely have commensurate running times. They start and end together and share nothing else. Interleaving them is possible and pointless: the audience shares a room but not a work.

At the next level, the realizations share a structure: act boundaries, intermissions, major turns. They may tell different stories, but the stories move in step, and viewers can sense that they are at comparable points.

At a stronger level, the realizations share a world: the same characters, places, and events, differing in how the events are presented. A comic and a dramatic realization of one plot are related this way.

At the strongest level, the realizations are generated from a common model of the story, a scene graph of places, characters, and events on a single timeline, and differ only in the choices made in rendering it: which camera, which line reading, which explanation, which outcome at a branch. Every difference is a declared transformation of a common source.

The book concentrates on the last two levels, because only they preserve what the Introduction called a single occasion. Strangers watching unrelated films in one room are co-present. Viewers watching two realizations of one event are participating in one differentiated work.

\section{Three kinds of distance}

Two realizations can be near in one sense and far in another, and it helps to separate three measures of how far apart they are.

\emph{Image distance} is how much of the time they show different pictures. Its complement is the coincidence fraction $c$ of \Cref{ch:capacity}, and it governs light and flash-rate costs. Two realizations that share most of their footage are cheap to multiplex.

\emph{Fact distance} is how much they differ in what they establish as having happened or being true: who did what, what was revealed, what the viewer is entitled to believe. It governs whether a viewer can move between them, which is the subject of \Cref{ch:switch-admissibility}.

\emph{Time distance} is how far their events drift apart on the common clock: whether a scene in one realization runs longer than the corresponding scene in another, so that the two fall out of step. It governs synchronization and the editing of \Cref{ch:editing}.

The three are analytically independent: each can in principle vary while the others are held fixed. \Cref{ch:divergence} shows that real transformations tend to couple them. A comic and a dramatic version of the same scene, shot separately, have large image distance and almost no fact distance. Two versions that use identical footage but differ in one subtitle revealing the murderer have almost no image distance and decisive fact distance. A beginner's version that pauses to explain a step has small fact distance but grows in time distance with every explanation. Any account of what a family costs, or what it allows, has to say which distance it is talking about.

\section{Primary and equal families}

A family may or may not have a privileged member. In a \emph{primary} family, one realization is the work in the conventional sense and the others are declared variants of it: an original with a children's adaptation, say, or a feature with an explanatory version for students. In an \emph{equal} family, no member is privileged; the work is the family.

The distinction is a matter of declaration, not of optics, and it has consequences later. It determines what an archive must preserve, what a critic reviews, and whether a viewer who chose a variant can be said to have seen the film. \Cref{ch:which-film} returns to it. Here it is enough to note that selective cinema does not force either answer. It makes the question unavoidable, where ordinary distribution lets it be deferred by releasing versions at different times to different audiences.

\section{A family is not a menu}

It is tempting to describe a variant family as a menu: a list of options from which each viewer picks. The description is misleading in two ways.

First, a menu's items are independent, and a family's are not. Each realization is constrained by the others through the shared scaffold. A scene in the comic realization cannot run long without the dramatic realization absorbing the difference somewhere, and a fact established in one cannot be contradicted in another without closing off the possibility of switching between them. The members of a family are written against each other.

Second, a menu offers a choice made once. A family offers a structure of choices across time. Whether a viewer can move from one realization to another, when, and in which direction depends on what each realization has shown so far. The family is therefore better understood as a set of paths through a common story, with a relation telling which paths can be joined to which at which moments. The next chapter defines that relation.
